\documentclass[10pt,a4paper]{amsart}
\usepackage{amsmath,amssymb,amsthm,graphicx}
\usepackage{verbatim}
\newenvironment{reference}{\comment}{\endcomment}
%\newenvironment{reference}{}{}
\newenvironment{slogan}{\comment}{\endcomment}
\newenvironment{history}{\comment}{\endcomment}

% For commutative diagrams we use Xy-pic
\usepackage[all]{xy}

% We use 2cell for 2-commutative diagrams.
\xyoption{2cell}
% Theorem environments.
%
\theoremstyle{plain}
\newtheorem{theorem}[subsection]{Theorem}
\newtheorem{proposition}[subsection]{Proposition}
\newtheorem{lemma}[subsection]{Lemma}

\theoremstyle{definition}
\newtheorem{definition}[subsection]{Definition}
\newtheorem{example}[subsection]{Example}
\newtheorem{exercise}[subsection]{Exercise}
\newtheorem{situation}[subsection]{Situation}

\theoremstyle{remark}
\newtheorem{remark}[subsection]{Remark}
\newtheorem{remarks}[subsection]{Remarks}

\numberwithin{equation}{subsection}

% Macros
%
\def\lim{\mathop{\mathrm{lim}}\nolimits}
\def\colim{\mathop{\mathrm{colim}}\nolimits}
\def\Spec{\mathop{\mathrm{Spec}}}
\def\Hom{\mathop{\mathrm{Hom}}\nolimits}
\def\Ext{\mathop{\mathrm{Ext}}\nolimits}
\def\SheafHom{\mathop{\mathcal{H}\!\mathit{om}}\nolimits}
\def\SheafExt{\mathop{\mathcal{E}\!\mathit{xt}}\nolimits}
\def\Sch{\mathit{Sch}}
\def\Mor{\mathop{\mathrm{Mor}}\nolimits}
\def\Ob{\mathop{\mathrm{Ob}}\nolimits}
\def\Sh{\mathop{\mathit{Sh}}\nolimits}
\def\NL{\mathop{N\!L}\nolimits}
\def\CH{\mathop{\mathrm{CH}}\nolimits}
\def\proetale{{pro\text{-}\acute{e}tale}}
\def\etale{{\acute{e}tale}}
\def\QCoh{\mathit{QCoh}}
\def\Ker{\mathop{\mathrm{Ker}}}
\def\Im{\mathop{\mathrm{Im}}}
\def\Coker{\mathop{\mathrm{Coker}}}
\def\Coim{\mathop{\mathrm{Coim}}}

% Boxtimes
%
\DeclareMathSymbol{\boxtimes}{\mathbin}{AMSa}{"02}

%
% Macros for moduli stacks/spaces
%
\def\QCohstack{\mathcal{QC}\!\mathit{oh}}
\def\Cohstack{\mathcal{C}\!\mathit{oh}}
\def\Spacesstack{\mathcal{S}\!\mathit{paces}}
\def\Quotfunctor{\mathrm{Quot}}
\def\Hilbfunctor{\mathrm{Hilb}}
\def\Curvesstack{\mathcal{C}\!\mathit{urves}}
\def\Polarizedstack{\mathcal{P}\!\mathit{olarized}}
\def\Complexesstack{\mathcal{C}\!\mathit{omplexes}}
% \Pic is the operator that assigns to X its picard group, usage \Pic(X)
% \Picardstack_{X/B} denotes the Picard stack of X over B
% \Picardfunctor_{X/B} denotes the Picard functor of X over B
\def\Pic{\mathop{\mathrm{Pic}}\nolimits}
\def\Picardstack{\mathcal{P}\!\mathit{ic}}
\def\Picardfunctor{\mathrm{Pic}}
\def\Deformationcategory{\mathcal{D}\!\mathit{ef}}
\usepackage{cleveref}
\crefname{originallemma}{original lemma}{original lemmas}
\Crefname{originallemma}{Original lemma}{Original lemmas}
\makeatletter
\newlabel{algebra-lemma-flatness-descends-more-general}{{0584}{source}{}{sourcepointer-1}{}}
\newlabel{algebra-lemma-flatness-descends-more-general@cref}{{[originallemma][0][]0584}{[1][0][]source}{}{}{}}
\makeatother
\crefname{originallemma}{original lemma}{original lemmas}
\Crefname{originallemma}{Original lemma}{Original lemmas}
\makeatletter
\newlabel{flat-lemma-flat-up-down-henselization}{{05VL}{source}{}{sourcepointer-2}{}}
\newlabel{flat-lemma-flat-up-down-henselization@cref}{{[originallemma][0][]05VL}{[1][0][]source}{}{}{}}
\makeatother
\crefname{originallemma}{original lemma}{original lemmas}
\Crefname{originallemma}{Original lemma}{Original lemmas}
\makeatletter
\newlabel{spaces-properties-lemma-describe-etale-local-ring}{{04KF}{source}{}{sourcepointer-3}{}}
\newlabel{spaces-properties-lemma-describe-etale-local-ring@cref}{{[originallemma][0][]04KF}{[1][0][]source}{}{}{}}
\makeatother
\crefname{originallemma}{original lemma}{original lemmas}
\Crefname{originallemma}{Original lemma}{Original lemmas}
\makeatletter
\newlabel{spaces-properties-lemma-geometric-lift-to-usual}{{05VN}{source}{}{sourcepointer-4}{}}
\newlabel{spaces-properties-lemma-geometric-lift-to-usual@cref}{{[originallemma][0][]05VN}{[1][0][]source}{}{}{}}
\makeatother
\title{Flatness of a morphism of algebraic spaces at a point equals flatness of the map on etale local rings}
\author{The Stacks Project Authors}
\date{stacks-project:spaces-morphisms.tex:local-snapshot-20261012; GNU FDL 1.2 or later}
\begin{document}
\maketitle
\setcounter{subsection}{1}
\noindent\textit{Editorial note.} All delivered mathematical text is the author's own, copied from the pinned Stacks Project chapter file spaces-morphisms.tex; no statement or proof step was rewritten, reordered or removed.\par\medskip
\noindent\textit{Editorial note.} The proof is the module-free companion of the flatness-at-a-point criterion; the strict henselization comparison is carried out explicitly inside the counted span.\par\medskip
\noindent\textit{Editorial note.} The four references of the counted spans are recorded as non-load-bearing navigation pointers; the argument that consumes them stays inside the counted proof span.\par\medskip
\noindent\textit{Editorial note.} Preamble: the three native excerpt spans of sources/stacks-spaces-morphisms/preamble.tex (lines 8-13, 14-18 and 154-217). The chapter file itself carries no preamble text.\par\medskip
\noindent\textit{Allowed standard prerequisites (editorial).}\begin{itemize}
\item Lifting geometric points along a surjective etale morphism. Properties of Spaces, Lemma spaces-properties-lemma-geometric-lift-to-usual, used at proof lines 6186-6191. Produces ubar over u with a o ubar = xbar.
\item Etale local rings as strict henselizations of the local rings of a presentation. Properties of Spaces, Lemma spaces-properties-lemma-describe-etale-local-ring, used at proof lines 6192-6200. Identifies the source and target of the ring map in the statement.
\item Flatness descends through the tensor description of the local rings. Algebra, Lemma algebra-lemma-flatness-descends-more-general, used at proof lines 6212-6215. Equates flatness of O\_\{U,u\} over O\_\{V,v\} with flatness of O\_\{X,xbar\} over O\_\{V,v\}.
\item The flat up-and-down property of henselization. More on Flatness, Lemma flat-lemma-flat-up-down-henselization, used at proof lines 6216-6216. Replaces O\_\{V,v\} by O\_\{Y,ybar\} and concludes.
\end{itemize}\medskip
\ifdefined\hypertarget\hypertarget{sourcepointer-1}{}\fi
\noindent\textit{Source reference.} Original lemma 0584: focused-seeds/AGP/sources/stacks-snapshot-20261011/tags.txt, source lines 6767--6767. Navigation pointer into the Stacks Project; the standard fact it supplies is applied inside the counted proof span and the argument resting on it stays there.\par\medskip
\ifdefined\hypertarget\hypertarget{sourcepointer-2}{}\fi
\noindent\textit{Source reference.} Original lemma 05VL: focused-seeds/AGP/sources/stacks-snapshot-20261011/tags.txt, source lines 7555--7555. Navigation pointer into the Stacks Project; the standard fact it supplies is applied inside the counted proof span and the argument resting on it stays there.\par\medskip
\ifdefined\hypertarget\hypertarget{sourcepointer-3}{}\fi
\noindent\textit{Source reference.} Original lemma 04KF: focused-seeds/AGP/sources/stacks-snapshot-20261011/tags.txt, source lines 5971--5971. Navigation pointer into the Stacks Project; the standard fact it supplies is applied inside the counted proof span and the argument resting on it stays there.\par\medskip
\ifdefined\hypertarget\hypertarget{sourcepointer-4}{}\fi
\noindent\textit{Source reference.} Original lemma 05VN: focused-seeds/AGP/sources/stacks-snapshot-20261011/tags.txt, source lines 7557--7557. Navigation pointer into the Stacks Project; the standard fact it supplies is applied inside the counted proof span and the argument resting on it stays there.\par\medskip
\begin{lemma}
\label{lemma-flat-at-point-etale-local-rings}
Let $S$ be a scheme.
Let $f : X \to Y$ be a morphism of algebraic spaces over $S$.
Let $\overline{x}$ be a geometric point of $X$ lying over the point
$x \in |X|$. Let $\overline{y} = f \circ \overline{x}$. The following
are equivalent
\begin{enumerate}
\item $f$ is flat at $x$, and
\item the map on \'etale local rings
$\mathcal{O}_{Y, \overline{y}} \to \mathcal{O}_{X, \overline{x}}$
is flat.
\end{enumerate}
\end{lemma}

\begin{proof}
Choose a commutative diagram
$$
\xymatrix{
U \ar[d]_a \ar[r]_h & V \ar[d]^b \\
X \ar[r]^f & Y
}
$$
where $U$ and $V$ are schemes, $a, b$ are \'etale, and
$u \in U$ mapping to $x$. We can find a geometric point
$\overline{u} : \Spec(k) \to U$ lying over $u$ with
$\overline{x} = a \circ \overline{u}$, see
Properties of Spaces, Lemma
\ref{spaces-properties-lemma-geometric-lift-to-usual}.
Set $\overline{v} = h \circ \overline{u}$ with image $v \in V$.
We know that
$$
\mathcal{O}_{X, \overline{x}} = \mathcal{O}_{U, u}^{sh}
\quad\text{and}\quad
\mathcal{O}_{Y, \overline{y}} = \mathcal{O}_{V, v}^{sh}
$$
see
Properties of Spaces, Lemma
\ref{spaces-properties-lemma-describe-etale-local-ring}.
We obtain a commutative diagram
$$
\xymatrix{
\mathcal{O}_{U, u} \ar[r] &
\mathcal{O}_{X, \overline{x}} \\
\mathcal{O}_{V, v} \ar[u] \ar[r] &
\mathcal{O}_{Y, \overline{y}} \ar[u]
}
$$
of local rings with flat horizontal arrows. We have to show that the
left vertical arrow is flat if and only if the right vertical arrow is.
Algebra, Lemma \ref{algebra-lemma-flatness-descends-more-general}
tells us $\mathcal{O}_{U, u}$ is flat over $\mathcal{O}_{V, v}$
if and only if $\mathcal{O}_{X, \overline{x}}$ is flat over
$\mathcal{O}_{V, v}$. Hence the result follows from
More on Flatness, Lemma \ref{flat-lemma-flat-up-down-henselization}.
\end{proof}

\end{document}
